%杨舒云的实验报告编辑界面，使用了Huanyu Shi，2019级的模板，杨舒云在此拜谢ORZ!

%!TEX program = xelatex
\documentclass[dvipsnames, svgnames,a4paper,11pt]{article}
\input{YsySettings} 
\usepackage{lipsum}
\usepackage{adjustbox}
\captionsetup[figure]{name=Figure} % 图片形式
\captionsetup[table]{name=Table} % 表格形式

\usepackage{enumitem}
\usepackage{tabularray}  %绘制表格时可以更加方便添加框线
\usepackage{xcolor} %添加更多文本颜色


%---------------------------------------------------------------------
%	正文
%---------------------------------------------------------------------



\begin{document}
	
	
	
	% 实验报告封面	
	
	% 顶栏
	\begin{table}
		\renewcommand\arraystretch{1.7}
		\begin{tabularx}{\textwidth}{
				|X|X|X|X
				|X|X|X|X|}
			\hline
			\multicolumn{2}{|c|}{预习报告}&\multicolumn{2}{|c|}{实验记录}&\multicolumn{2}{|c|}{分析讨论}&\multicolumn{2}{|c|}{总成绩}\\
			\hline
			\LARGE25 & & \LARGE30 & & \LARGE25 & & \LARGE80 & \\
			\hline
		\end{tabularx}
	\end{table}
	% ---
	
	% 信息栏
	\begin{table}
		\renewcommand\arraystretch{1.7}
		\begin{tabularx}{\textwidth}{|X|X|X|X|}
			\hline
			年级、专业： & 2022级 物理学 &组号： & 实验班2\\
			\hline
			姓名： & 杨舒云  & 学号： & 22344020\\
			\hline
			实验时间： & 2024/4/25 & 教师签名： & \\
			\hline
		\end{tabularx}
	\end{table}
	% ---
	
	% 大标题
	\begin{center}
		\LARGE Lab2-7 \quad 热辐射实验\\The Experiments on Thermal Radiation
	\end{center}
	% ---
	
	% 注意事项
	
	% 基本
	\textbf{【实验报告注意事项】}
	\begin{enumerate}
		\item 实验报告由三部分组成：
		\begin{enumerate}
			\item 预习报告：课前认真研读实验讲义，弄清实验原理；实验所需的仪器设备、用具及其使用、完成课前预习思考题；了解实验需要测量的物理量，并根据要求提前准备实验记录表格（可以参考实验报告模板，可以打印）。\textcolor{red}{\textbf{（25分）}}
			\item 实验记录：认真、客观记录实验条件、实验过程中的现象以及数据。实验记录请用珠笔或者钢笔书写并签名（\textcolor{red}{\textbf{用铅笔记录的被认为无效}}）。\textcolor{red}{\textbf{保持原始记录，包括写错删除部分，如因误记需要修改记录，必须按规范修改。}}（不得输入电脑打印，但可扫描手记后打印扫描件）；离开前请实验教师检查记录并签名。\textcolor{red}{\textbf{（30分）}}
			\item 数据处理及分析讨论：处理实验原始数据（学习仪器使用类型的实验除外），对数据的可靠性和合理性进行分析；按规范呈现数据和结果（图、表），包括数据、图表按顺序编号及其引用；分析物理现象（含回答实验思考题，写出问题思考过程，必要时按规范引用数据）；最后得出结论。\textcolor{red}{\textbf{（25分）}}
		\end{enumerate}
		\textbf{实验报告就是将预习报告、实验记录、和数据处理与分析合起来，加上本页封面。\textcolor{red}{（80分）}}
		\item 每次完成实验后的一周内交\textbf{实验报告}（特殊情况不能超过两周）。
	\end{enumerate}
	
	% 安全
	\textbf{【实验安全注意事项】}	
	\begin{enumerate}
		\item 实验过程中，禁止触摸辐射体表面，一方面是避免在高温时烫伤；另一方面避免污染表面，影响发射系数。
		\item 测量不同辐射表面对辐射强度影响时，辐射温度不要设置太高，更换辐射体时，应带手套。
		\item 实验过程中，计算机在采集数据时不要触摸测试架，以免造成对传感器的干扰。
		\item 在使用控温程序时，要留意程序是否正常运行，若运行有问题，请及时关闭可编程直流电源的输出，防止辐射器温度过高而损坏。停止运行程序时，要用程序的Stop output按钮， 不要用菜单键的红色强制终止按钮。
		\item 辐射传感器上方有一块金属挡板。在测量时将挡板移开；非测量时关上，避免不必要的触碰污染传感器表面。
	\end{enumerate}
	
	% ---
	
	% 特别鸣谢
	\textbf{【特别鸣谢及模板说明】}	
	
	感谢2019级学长石寰宇为本实验报告提供\LaTeX 模板。\textcolor{red}{\textbf{由于原实验报告模板缺少实验编号，为方便在电脑上整理，故添加自命名编号Lab2-}}
	% ---
	
	
	
	% 目录
	\clearpage
	\tableofcontents
	\clearpage
	% ---
	
	
	
	% 预习报告	
	
	% 小标题
	\setcounter{section}{0}
	\section{Lab2-7 热辐射实验（The Experiments on Thermal Radiation）\\ \quad\heiti 预习报告（Pre-Report）}
	% ---
	
	% 实验目的
	\subsection{实验目的（Purpose）}
	\begin{enumerate}
		\item 认识热辐射现象及其本质（普遍存在的一种能量转换与传递的形式）。
		\item 认识影响热辐射强度的各种因素及其与热辐射强度的定量关系，包括：
		
		\begin{enumerate}
			\item 辐射体表面温度（验证Stefan-Boltzmann定律）；
			\item 辐射距离（面源辐射修正）；
			\item 表面的发射系数（拓展）；
			\item 验证上述因素与辐射强度的定量关系是否符合黑体辐射定律？
		\end{enumerate}
		
		\item 了解热辐射传感器（SMTIR9902）原理和结构、使用（含校正）方法。
		\item 学习应用LabView管理由具有NI通信协议的非NI专业仪器（数字多用表）、设备（程控电源）构成的实验系统。
	\end{enumerate}
	% ---
	
	% 仪器用具
	\subsection{仪器用具（Instruments）}
	\begin{table}[htbp]
		\centering
		\renewcommand\arraystretch{1.6}
		% \setlength{\tabcolsep}{10mm}
		\begin{tabular}{p{0.05\textwidth}|p{0.20\textwidth}|p{0.05\textwidth}|p{0.5\textwidth}}
			\hline
			编号& 仪器用具名称 & 数量 &  主要参数（型号，测量范围，测量精度等） \\
			\hline
			1&黑体辐射与红外测量装置 	&1 	& DHRH-B：含带标尺（60cm）位移导轨、辐射器（三种辐射面：黑面、粗面、光滑金属面）、热辐射传感器（SMTIR9902），\\
			
			2&数字多用表 & 2	& RIGOL DM3058E \\
			
			3&程控电源 	& 1 & RIGOL DP831\\
			
			4&计算机 &1 &已安装 LabView 和控温软件 \\
			\hline
		\end{tabular}
	\end{table}
	% ---
	
	% 原理概述
	\subsection{原理概述（Principle）}
	
	%
	\subsubsection{热辐射的产生}
	\begin{enumerate}
		\item 辐射本领和吸收本领的关系: G.R.Kirchhoff的热辐射定律表明，一个物体对于特定波长和温度的热辐射本领与其吸收本领成正比，这个比值仅与频率和温度相关。
		
		\item 黑体辐射: 任何物体都能辐射和吸收能量。一个理想化的黑体是指可以吸收所接受到的所有辐射能量而不反射和透过任何辐射的物体。黑体辐射的能量与物体的温度有关，其辐射特性只由温度决定。
		
		Planck公式描述了一个完美黑体在温度为 \( T \) 时，单位面积上单位时间内在单位立体角内发射的电磁辐射能量密度（辐射强度）。公式表达如下：
		\[
		B_\nu(T) = \frac{2h\nu^3}{c^2} \frac{1}{e^{\frac{h\nu}{kT}} - 1}
		\]
		其中：\( B_\nu(T) \) 是频率 \( \nu \) 下的辐射强度；\( h \) 是Planck常数；\( \nu \) 是辐射的频率；\( c \) 是光速；\( k \) 是Boltzmann常数；\( T \) 是黑体的绝对温度。
		
		\item Stefan-Boltzmann定律: 描述了黑体单位面积的总辐射强度与其绝对温度的四次方成正比的关系，公式为 \( R = \sigma T^4 \) 其中 \( \sigma \) 是Stefan-Boltzmann常数。
		
		推导如下：
		
		黑体在所有频率上的辐射强度总和可以通过对Planck公式进行积分得到：
		
		\[
		M = \int_0^\infty B_\nu(T) d\nu
		\]
		
		通过变量替换和积分，可以得到：
		
		\[
		M = \frac{2\pi^5 k^4}{15h^3c^2} T^4
		\]
		
		Stefan-Boltzmann常数 \( \sigma \) 定义为 \( \frac{2\pi^5 k^4}{15h^3c^2} \)。
		
		\item Wien位移定律: 描述了黑体辐射光谱中强度最大的波长与绝对温度成反比的关系，表明随着温度的增加，最大辐射强度的波长向短波方向移动。
		
		Wien位移定律表明黑体辐射强度峰值对应的波长与温度成反比，公式为：
		
		\[
		\lambda_{\text{max}} T = b
		\]
		
		其中 \( b \) 是Wien位移常数。
		
		推导如下：
		
		求Planck公式对波长的极大值：
		
		\[
		\frac{d}{d\lambda} B_\lambda(T) = 0
		\]
		
		将Planck公式由频率形式转为波长形式，并对其求导。使用相关的代数操作和近似方法，可以找到使得导数为零的波长。
		
		经过数学计算后，可以得出 \( \lambda_{\text{max}} T \approx 2.898 \times 10^{-3} \) 米·开尔文。
		
		\item 能量转换与传递:通过电加热器，电能转换成热能，热能由辐射体的表面以电磁波形式辐射出去。辐射能量的大小和分布受到物体的温度、表面性质（如黑度、材质）和形状等多因素影响。
	\end{enumerate}
	
	%
	\subsubsection{热辐射的测量}
	\begin{enumerate}
		\item 根据Stefan-Boltzmann定律， 辐射体的辐射本领与其绝对温度的四次方成正比。实验中，通过电加热器将电能转换为热能，通过辐射面发射，并由热辐射传感器捕获。
		
		\item 热辐射传感器基于热电堆原理，使用多个串联的热电偶。当辐射体发射的热辐射到达热电堆时，吸热面温度升高产生温度差，进而产生热电势。传感器输出的电压信号与接收到的辐射强度成正比，从而测量热辐射。
		
		具体来说，SMTIR9902系列传感器通过热电堆（一系列串联的热电偶）测量热辐射。当辐射体向传感器的吸热面辐射热量时，吸热面的温度上升，形成与传感器冷端（通常与环境温度接近或通过特定冷却措施控制）的温差。这个温差产生一个电势差，即热电势，该电势被转换为电压输出，电压的大小与接收的辐射强度成正比。
		\[
		(\dot{Q} = \dot{q} A = \kappa\Delta 𝑇 \propto \Delta V)
		\]
		因此，SMTIR9902的输出信号正比于当地从辐射方向进入的辐射通量或能流密度，或接收功率$\Delta V=SQ$， 其中S为探测灵敏度；从说明书得到，$ S=110\pm20(V/W)@500K$。传感器接收面积已确定$(A= 0.5{mm}^2)$，从而可以计算当地辐射通量和对辐射源的立体角，再进一步计算辐射源辐射强度（公式以 (W/sr) 为单位）
		\[
		I_E= d\varphi /d\Omega=\dot{q}d^2
		\]
		
		\quad
		
		\textbf{【讨论：关于传感器测量功能的拓展】}
		\begin{enumerate}
			\item 该传感器是否适合测量高温热辐射：SMTIR9902系列传感器设计用于低温至中温范围的热辐射测量。这是因为它的动态范围和设计主要针对这一温区的辐射特性。对于高温热辐射，传感器可能因为接收到的辐射能量过大而受损，或其读数超出设计的测量范围，从而导致数据不准确或设备损坏。高温热辐射通常伴随更高的能量输出和更宽的光谱范围。因此，如果要测量高温热辐射，建议使用专为高温设计的传感器，这些传感器具有更高的耐温性能和更广的响应范围。
			
			\item 当辐射体的温度低于环境温度（室温，即传感器温度）时，辐射动态将发生变化：辐射体成为热量接收者，此时，辐射体不再是主要的热辐射源，而是可能从周围环境（包括传感器）接收热辐射。传感器的吸热面可能会失去热量，如果传感器的吸热面比辐射体温度高，吸热面可能会向较冷的辐射体失去热量，导致吸热面的温度降低。因为热电堆是基于温差工作的，当吸热面的温度因为辐射体的较低温度而下降时，与冷端的温差减少，从而导致热电势减小，输出电压相应减少。在某些情况下，如果辐射体温度远低于传感器温度，理论上可能导致传感器输出信号的方向发生变化（即从输出正电压变为输出负电压），这取决于传感器设计和接线方式。
			
			\item 根据讲义中的内容，尽管按照物理定律，处于室温的物体也会发出热辐射，但是当物体的温度和环境温度相同时，物体的辐射和吸收达到平衡。此时，传感器可能测得的辐射强度与环境的背景辐射强度相等，因此在没有温差的情况下，辐射体和传感器之间没有净能量交换，导致传感器输出的电压信号为零。实际上，这是因为传感器主要测量的是温差引起的辐射强度变化，而非绝对辐射强度。
			
			\item 要直接测量处于室温的物体的辐射强度，理想情况下，传感器外壳（环境）的温度应该低于被测物体的温度，以便产生足够的温差，从而使传感器能够检测到由被测物体向环境（或传感器）发出的净辐射。通过增加环境和物体之间的温差，可以增加传感器的输出信号，使其能够测量到来自物体的辐射。
		\end{enumerate}
		
		\item 测量误差来源及处理手段
		\begin{enumerate}
			\item 消除人体热辐射对SMTIR9902系列传感器的影响：人体作为一种热源，其热辐射可能会影响传感器的读数，特别是在测量较低温度的辐射时。为消除这种影响，可以采用物理隔离的方法，如在传感器与人体之间设置屏障或使用反射材料来遮挡人体辐射。此外，确保实验环境足够大，使传感器远离操作者和其他人员，也是减少人体热辐射影响的有效方式。另外，可以通过校准来补偿人体热辐射的影响，例如在实验开始前记录传感器在无其他热源干扰时的基线数据，并在实验数据分析时考虑这一基线影响。
			
			\item 日光灯对SMTIR9902系列传感器的影响：通常情况下，日光灯主要发出可见光和少量紫外线，对于专门测量红外辐射的传感器如SMTIR9902系列，其对可见光的响应非常低。该传感器通常配备有过滤器，可以过滤掉大部分可见光，从而减少日光灯等光源的影响。但是，如果实验环境中存在强烈的光照，仍建议采用遮光措施，如使用遮光帘或将实验设备放置在无直接光照的环境中，以避免任何潜在的光学干扰。
			
			\item SMTIR9902SIL自带透镜，减少接收角从而减少非探测物辐射的影响。
			
			\item 环境温度变化的数学扣除方法
			
			如果环境温度与室温物体的温度接近，且有变化，可以通过数学方法来扣除环境温度的影响。
			这可以通过校准传感器来实现，具体方法包括：
			\begin{itemize}
				\item 记录基线数据：在无其他辐射源的条件下，记录环境温度对应的传感器输出，作为基线。
				\item 建立数学模型：建立传感器输出与环境温度之间的关系模型。这可以通过在不同已知环境温度下测量传感器输出，然后使用回归分析或其他统计方法来建模。
				\item 温度补偿：使用上述模型来调整实际测量值，从而补偿由环境温度变化引起的传感器输出变化。
			\end{itemize}
			这种方法需要精确的温度传感器（如实验中提到的内置Ni1000热敏电阻），以及精确记录和调整环境温度变化。通过这种方法，可以有效地从传感器的总输出中分离出由环境温度变化引起的部分，从而更准确地测量物体的辐射强度。
	
		\end{enumerate}
	\end{enumerate}
	
	%
	\subsubsection{实验方案设计}
	\begin{enumerate}
		\item 实验设备和配置：实验装置包括一个可调节距离的辐射源和热辐射传感器，两者通过计算机控制。
		
		\begin{itemize}
			\item 辐射体：通过电加热方式控制其温度，并使用PT1000温度传感器实时监控辐射体表面温度。辐射体温度可以高达120℃，因此操作时需要避免直接接触。
			
			\quad
			
			\textbf{【讨论：关于辐射体温度与加热性质】}
			\begin{itemize}
				\item 辐射体的加热功率与辐射体温度之间的关系：
				
				Stefan-Boltzmann定律是理解辐射体加热功率与其温度关系的关键。该定律指出辐射体的总辐射能量与其绝对温度的四次方成正比，表达式为 \( R = \sigma T^4 \)，其中 \( R \) 是辐射功率，\( T \) 是绝对温度，而 \( \sigma \) 是Stefan-Boltzmann常数。
				
				在实验设置中，辐射体被加热至特定温度，加热功率（电功率）主要用于维持辐射体温度对抗周围环境的热损失。随着辐射体温度的增加，为保持该温度所需的功率也随之增加，尤其是当温度高于环境温度时，由于热损失（主要是通过辐射和对流）增加。
				
				\item 加热功率对传感器信号的影响：
				
				辐射传感器（如SMTIR9902系列）的信号输出主要依赖于其接收到的辐射能量，这与辐射体的温度密切相关。因此，当辐射体的温度因加热功率的增加而升高时，其辐射强度增加，导致传感器检测到更多的热能，并输出更高的电压信号。
				
				传感器的输出信号与辐射体的温度基本上呈正比关系。这是因为传感器设计为测量辐射体表面的辐射强度，而这一强度根据Stefan-Boltzmann定律与温度的四次方成正比。
				
				\item 物理上我们可以考虑以下两个因素：
				
				能量守恒：加热功率提供了维持或增加辐射体温度所需的能量，这些能量通过辐射、对流和传导的形式传递到环境中。辐射形式的能量转换直接影响了传感器的检测。
				
				热辐射的物理性质：辐射体的温度提高导致其发射更多的红外辐射，这增加了辐射传感器所检测到的辐射强度，并转化为更强的电信号输出。
			\end{itemize}
			
			\item 辐射测量：热辐射传感器固定在可调节的距离上，以测量来自辐射体的辐射强度。（实验中可以调节辐射体与传感器之间的距离，以研究距离变化对辐射测量的影响。）辐射距离和角度的调整是实验中重要的变量，可以影响测量结果的精确性。
			
			\item 实验中应用金属挡板遮光，以避免环境光干扰传感器（详见前讨论）。实验系统中，热辐射传感器安装在可移动的金属铝臂上，保证传感器外壳与铝臂良好的热接触，而传感器的位置可通过导轨上方的手摇臂连续调节；辐射传感器上方提供了一块金属挡板。在测量时可将挡板移开；非测量时关上，避免不必要的触碰污染传感器表面，给测量造成误差。
			
			\quad
			
			\textbf{【讨论：关于防热辐射实验内容】}
			
				Beer定律：这是一个描述光（或其他电磁辐射）通过吸收介质传播时强度衰减的定律。它表明透射的光强度 \(I\) 与介质的厚度 \(d\) 成指数关系，可以表示为：
				\[
				I = I_0 e^{-\mu d}
				\]
				其中，\(I_0\) 是未经材料吸收前的初始光强度，\(\mu\) 是材料的吸收系数，\(d\) 是材料的厚度。
				
				如果对上述等式取自然对数，我们得到：
				\[
				\ln(I) = \ln(I_0) - \mu d
				\]
				这表明 \(\ln(I)\) 与材料的厚度 \(d\) 是线性关系。
				
				因此考虑以下实验方案：
				\begin{enumerate}
					\item \textbf{实验设置：}
					\begin{itemize}
						\item 使用一系列不同厚度的防热辐射材料。
						\item 设置辐射体（如黑体辐射器）在一个固定温度。
						\item 使用热辐射传感器（如SMTIR9902系列）来测量透过这些材料的辐射强度。
					\end{itemize}
					
					\item \textbf{数据采集：}
					\begin{itemize}
						\item 对每种不同厚度的材料，记录辐射传感器的读数。
						\item 可以使用更精确的仪器如红外光谱仪来获得更详细的辐射强度数据。
					\end{itemize}
					
					\item \textbf{数据分析：}
					\begin{itemize}
						\item 将测得的辐射强度 \(I\) 与对应的材料厚度 \(d\) 记录下来。
						\item 计算每个数据点的 \(\ln(I)\)，并与厚度 \(d\) 作图。
						\item 使用线性回归分析来确定直线的斜率，该斜率即为 \(-\mu\)。
					\end{itemize}
					
					\item \textbf{实验验证：}
					\begin{itemize}
						\item 通过实验数据，验证辐射强度与材料厚度的指数衰减关系是否成立。
						\item 通过斜率 \(\mu\)，可以进一步了解材料的防辐射性能，较大的 \(\mu\) 值意味着更强的辐射吸收能力。
					\end{itemize}
				\end{enumerate}
				
		\end{itemize}		
		\item 数据采集与控制：实验中使用LabVIEW软件进行数据采集和实验控制。通过计算机与实验设备（如数字多用表和程控电源）的通信，实时采集辐射体的温度和传感器的电压数据，自动记录和处理实验数据。
		
		\item 实验步骤
		
		\begin{enumerate}
			\item 设置辐射体的温度并等待其稳定。
			\item 通过LabVIEW程序控制实验，包括温度的设置和辐射传感器的读取。
			\item 改变实验条件如辐射距离和表面类型，观察对热辐射测量的影响。
		\end{enumerate}
	\end{enumerate}
	
	
	% ---
	
	
	
	% 实验前思考题
	\subsection{实验前思考题（Questions）}
	
	% 思考题1
	\begin{question}
		人体热辐射会对SMTIR9902系列传感器读数产生影响（自己可验证），如何消
		除这种影响？日光灯是否有影响呢？这种传感器是否适合测量高温热辐射，
		为什么？
	\end{question}
	详见\textbf{原理概述}部分。
	
	% 思考题2
	\begin{question}
		如果对辐射体制冷，使辐射表面温度低于室温（传感器温度） ，辐射传感器
		输出信号会如何变化？
	\end{question}
	详见\textbf{原理概述}部分。
	
	% 思考题3
	\begin{question}
		按辐射定律，处于室温的物体也有辐射，为何辐射传感器的输出信号为零？
		如果要直接测量该物体的辐射强度（辐射传感器输出信号正比于辐射强
		度），环境（传感器外壳）温度应该多高？ 如果环境（传感器外壳）温度维
		持在室温附近但温度有变化， 是否可以通过数学方法扣除（非绝对零度的）
		环境温度的影响？（对给定辐射源温度， 传感器外壳温度与辐射传感器读数
		之间是什么关系？ 提示，热辐射传感器内置Ni1000可以测量外壳的温度，其
		分度表见附录2）
	\end{question}
	详见\textbf{原理概述}部分。
	
	% 思考题4
	\begin{question}
		辐射体的加热功率与辐射体温度之间呈何关系？与辐射传感器的信号值之间
		呈何关系？为什么？
	\end{question}
	详见\textbf{原理概述}部分。
	
	% 思考题5
	\begin{question}
		（选） （对防热辐射实验内容）透过防热辐射材料后的辐照度随材料的厚度
		怎么变化？请推导并验证。
	\end{question}
	详见\textbf{原理概述}部分。
	
	% ---
	
	
	
	% 实验记录	
	\clearpage
	
	% 顶栏
	\begin{table}
		\renewcommand\arraystretch{1.7}
		\centering
		\begin{tabularx}{\textwidth}{|X|X|X|X|}
			\hline
			专业： & 物理学 & 年级： & 2022级 \\
			\hline
			姓名： & 杨舒云 & 学号： & 22344020\\
			\hline
			室温： & 26\degree C & 实验地点： & A515 \\
			\hline
			学生签名：& 见\textbf{附件} & 评分： &\\
			\hline
			实验时间：& 2024/4/25 & 教师签名：&\\
			\hline
		\end{tabularx}
	\end{table}
	% ---
	
	% 小标题
	\section{Lab2-7 热辐射实验（The Experiments on Thermal Radiation）\\ \quad\heiti 实验记录（Records）}
	% ---
	
	% 实验过程记录
	\subsection{实验内容、步骤与结果（Content, Procedures \& Results）}
	
	%
	\subsubsection{操作步骤记录（Operations）}
	\begin{enumerate}
		\item 实验一 实验系统搭建
		\begin{enumerate}
			\item 辐射体中的加热模块中有两个$10\Omega$电阻，他们两个可能是独立的，也可能是共地的，我们需要用万用表欧姆档判断哪两条线之间对应一个电阻；判断完后，将两个电阻并联接入电源，使用并联的目的是增大功率，加快加热速度。
			\item 辐射体中的温度传感器模块同样有两个$PT1000$电阻，同样使用万用表欧姆档判别，随后只需将一个$PT1000$电阻接入一个万用表。
			\item 辐射传感器部分有一个热电堆和一个$NI1000$电阻，在本次实验中使用热电堆最为辐射传感器，将其接入一个万用表。
			\item 所有的接线完成后，将辐射体和辐射传感器安装在实验架上。
			\item 检查数据线是否与电脑程序正常连接。
			\item 所有准备工作完成后，开始下面的实验步骤。
		\end{enumerate}
		
		\item 实验二 测量物体的辐射面温度对辐射强度大小的影响
		\begin{enumerate}
			\item 选择辐射距离为\textbf{140mm}，将辐射传感器调至该距离；
			
			\item 设计数个温度控制点；由于120\degree C的断电保护控制，建议设计的控温温度设置值不高于118\textdegree C。记录每个控温点（稳定后）所需的电功率以及辐射传感器输出电压；
			
			\item 分析数据，探索辐射强度与辐射表面温度之间的关系，并绘制温度-辐射强度曲线图；
			
			\item 分析数据，探索电源加热功率与辐射表面温度之间的关系。
		\end{enumerate}
		
		\item 实验三 测量在不同辐射距离在d的辐射传感器输出P 
		\begin{enumerate}
			\item 将控温度控制目标设置温度120\degree C，依然选择用黑面辐射体进行实验。
			
			\item 设计数个位置测量点。
			
			\item 移动红外传感器的位置，每移动一定的距离后，记录测得对应辐射强度，绘制$P-d^{2}$图。			
		\end{enumerate}
		
		\item 实验四 测量不同物体表面的发射系数
		
		将热辐射传感器移开，更换所需辐射面的辐射体。分别正确连接辐射器PT1000和陶瓷加热器的连线到万用表和可编程直流电源上，控温表设置在\textbf{100℃}，待温度控制好后，将热辐射传感器移至靠近辐射体\textbf{140mm}处，测量并记录该辐射传感器读数（实验时，保证热辐射传感器与待测辐射面距离相同，便于分析和比较）。
	\end{enumerate}	
	
	%
	\subsubsection{实验结果展示（Display）}
	\begin{enumerate}
		\item 实验系统搭建与软件控制界面
		
		搭建好的实验系统实物图如\cref{fig:fig1}所示。测量过程中的软件控制界面可以观察到如\cref{fig:fig2}所示。
		
		\begin{figure}[htbp]
			\centering
			\includegraphics[width=0.5\textwidth]{Lab2_7Gra1.png}
			\caption{实验系统实物图}
			\label{fig:fig1}
		\end{figure}
		
		\begin{figure}[htbp]
			\centering
			\subfloat[]{
				\includegraphics[width=0.4\textwidth]{Lab2_7Gra2.png}
			}
			\subfloat[]{
				\includegraphics[width=0.4\textwidth]{Lab2_7Gra3.png}
			}
			\caption{软件控制界面}
			\label{fig:fig2}			
		\end{figure}
		
		\item 测量物体的辐射面温度对辐射强度大小的影响
		
		参数说明：实验中将辐射传感器固定在 140.0mm 的位置。
		
		实验过程中记录的数据如\cref{tab:tab1}所示。
		
		\begin{table}[h]
			\centering
			\caption{实验数据：温度与传感器输出电压、电源输出功率的关系}
			\label{tab:tab1}
			\begin{tabular}{|c|c|c|}
				\hline
				设定温度/摄氏度 & 传感器输出电压/微伏 & 电源输出功率/瓦 \\
				\hline
				40 & 79 & 0.50 \\
				50 & 105 & 0.90 \\
				60 & 131 & 1.30 \\
				70 & 167 & 1.70 \\
				80 & 200 & 1.95 \\
				90 & 240 & 2.25 \\
				100 & 283 & 2.80 \\
				110 & 330 & 3.10 \\
				120 & 369 & 4.00 \\
				\hline
			\end{tabular}
		\end{table}
		
		实际上，由于软件本身也记录了数据，所以在后面的数据处理过程中我使用的是软件记录的数据，用法是\textbf{取一系列数据在一段时间内的平均值}，详见数据分析部分。
		
		\item 测量在不同辐射距离在d的辐射传感器输出P
		
		参数说明：在上个试验结束后，辐射体温度已经上升到 120℃，此时保持温度不变。
		
		实验过程中记录的数据如\cref{tab:tab2}所示。
		
		\begin{table}[h]
			\centering
			\caption{实验数据：距离与传感器输出电压的关系}
			\label{tab:tab2}
			\begin{tabular}{|c|c|}
				\hline
				距离/毫米 & 传感器输出电压/微伏 \\
				\hline
				120 & 512 \\
				130 & 426 \\
				140 & 366 \\
				150 & 315 \\
				160 & 275 \\
				170 & 240 \\
				180 & 213 \\
				190 & 200 \\
				200 & 178 \\
				210 & 166 \\
				220 & 153 \\
				230 & 143 \\
				240 & 139 \\
				250 & 135 \\
				\hline
			\end{tabular}
		\end{table}		
		
		实际上，由于软件本身也记录了数据，所以在后面的数据处理过程中我使用的是软件记录的数据，用法是\textbf{取一系列数据在一段时间内的平均值}，详见数据分析部分。
		
		\item 测量不同物体表面的发射系数
		
		参数说明：我们将不同的辐射体升高至相同的温度（保持不同的辐射体均在100℃），将辐射传感器放置在相同的位置上（实验中将辐射传感器固定在 140.0mm 的位置），通过比较辐射传感器的输出电压值的不同，判断不同辐射体的表面发射系数大小关系。
		
		实验过程中记录的数据如\cref{tab:tab3}所示。
		
		\begin{table}[h]
			\centering
			\caption{实验数据：不同类型表面的传感器输出电压}
			\label{tab:tab3}
			\begin{tabular}{|c|c|c|c|}
				\hline
				类型 & 光面 & 粗糙面 & 黑面 \\
				\hline
				传感器输出电压/微伏 & 285 & 313 & 1030 \\
				\hline
			\end{tabular}
		\end{table}
		
	\end{enumerate}
	
	% ---
	
	% 原始数据
	%\clearpage
	\subsection{原始数据记录（Original Data）}
	实验记录本上的原始数据见\cref{fig:figA1}与\cref{fig:figA2}（签字）。
	
	\begin{figure}[htbp]
		\centering
		\includegraphics[width=0.6\textwidth]{Lab2_7GraA1.jpg}
		\caption{原始数据1}
		\label{fig:figA1}
	\end{figure}
	
	\begin{figure}[htbp]
		\centering
		\includegraphics[width=0.5\textwidth]{Lab2_7GraA2.jpg}
		\caption{原始数据2}
		\label{fig:figA2}
	\end{figure}
	
	%实验台桌面整理见%\textbf{附件}部分(\cref{})。
	
	%其它原始数据见%\cref{}。
	% ---
	
	% 问题记录
	\subsection{实验过程中遇到的问题记录（Difficulties）}
	\begin{enumerate}
		\item 需要检查确定三台实验仪器与电脑之间正常连接。
		\item 在正式开始升温测量前，一定要确保接线正确。如果接线错误，可能会导致程序错误判断辐射体当前温度，导致升温过度损坏实验设备。我的实验仪器上便是黑与黑一组，红与红一组。
	\end{enumerate}
	% ---
	
	
	
	% 分析与讨论	
	\clearpage
	
	% 顶栏
	\begin{table}
		\renewcommand\arraystretch{1.7}
		\begin{tabularx}{\textwidth}{|X|X|X|X|}
			\hline
			专业：& 物理学 &年级：& 2022级\\
			\hline
			姓名： & 杨舒云 & 学号：& 22344020\\
			\hline
			日期：& 2024/4/30 & 评分： &\\
			\hline
		\end{tabularx}
	\end{table}
	% ---
	
	% 小标题
	\section{Lab2-7 热辐射实验（The Experiments on Thermal Radiation）\\ \quad\heiti 分析与讨论（Analysis \& Discussion）}
	% ---
	
	% 数据处理
	\subsection{实验数据分析（Analysis）}
	
	%
	\subsubsection{根据测量物体的辐射面温度对辐射强度大小的影响，分析讨论$P-T$关系（Analysis）}
	\begin{enumerate}
		\item 根据Stefan-Boltzmann定律（详见\textbf{原理概述}部分，考虑拟合函数$P=k(t+273.15+t_{bias})^4+P_{bias}$，其中t是测得的摄氏温度，k是包含系数$\sigma$的某个比例系数，而从代码运行的结果来看，需要加入两个偏置项。
		当辐射体的温度恒定时，电源的稳定输出功率值就等于辐射体向外辐射能量的功率，其应当满足$P\propto T^4$的正比关系。因此采用同样的拟合方式。然而，拟合结果表面，上述假定是不成立的，从曲线拟合结果可以看出电源的稳定输出功率值反而与温度呈现出一定的线性关系，因此重新考虑线性拟合。
		
		\item 利用处理后的数据（\cref{tab:tab4})，得到了如\cref{fig:fig4}、\cref{fig:fig5}和\cref{fig:fig6}所示的拟合结果。
		
		\begin{table}[h]
			\centering
			\caption{数据分析：表面温度与辐射传感器电压、电源功率的关系}
			\label{tab:tab4}
			\begin{tabular}{|c|c|c|}
				\hline
				表面温度 & 辐射传感器电压/mV & 电源功率/W \\
				\hline
				40.000 & 79.000 & 0.500 \\
				50.000 & 105.000 & 0.900 \\
				60.000 & 131.000 & 1.300 \\
				69.987 & 166.015 & 1.681 \\
				79.999 & 201.094 & 1.989 \\
				89.999 & 239.214 & 2.370 \\
				99.991 & 280.939 & 2.805 \\
				109.997 & 322.964 & 3.135 \\
				119.998 & 371.444 & 3.587 \\
				\hline
			\end{tabular}
		\end{table}
		
		
		\begin{figure}[htbp]
			\centering
			\includegraphics[width=0.6\textwidth]{Lab2_7Gra4.png}
			\caption{数据分析：辐射面温度对辐射强度大小的影响}
			\label{fig:fig4}
		\end{figure}
		
		\begin{figure}[htbp]
			\centering
			\includegraphics[width=0.6\textwidth]{Lab2_7Gra5.png}
			\caption{数据分析：电源功率的曲线拟合（失败）}
			\label{fig:fig5}
		\end{figure}
		
		\begin{figure}[htbp]
			\centering
			\includegraphics[width=0.6\textwidth]{Lab2_7Gra6.jpg}
			\caption{数据分析：电源功率的直线拟合}
			\label{fig:fig6}
		\end{figure}
		
		从相关系数可以看出，上述拟合效果较好，说明实验数据与物理规律符合程度较好。
		
		\item 接下来讨论分析这些拟合参数。
		
		辐射面温度对辐射强度大小的影响：从拟合参数可以看出，环境带来的背景辐射约128$\mu V$。至于温度项的偏置为什么不为0，我猜测有两种可能，一种是传感器电压与接收到的辐射功率是线性关系但不成正比；另一种可能是PT1000并不是严格满足理论上的电阻温度关系，由它引入这个偏置。
		
		而对于电源功率为什么与温度成线性关系，目前我并没有很好的想法。
		
	\end{enumerate}
	
	%
	\subsubsection{根据测量在不同辐射距离在$d$的辐射传感器输出$P$，分析讨论$P-d$关系（Analysis）}
	\begin{enumerate}
		\item 假设我们考虑辐射体发出的辐射是以球面波的形式传播出去的，则辐射传感器的电压值与距离辐射体的距离应当满足平方反比关系，考虑拟合函数$P=\frac{P_0}{(d-d_{bias})^{2}}+P_{bias}$，同样根据代码运行结果来看，需要加入两个偏置。
		
		\item 利用处理过后的数据（\cref{tab:tab5}）得到拟合后的结果如\cref{fig:fig7}所示。
		
		\begin{table}[h]
			\centering
			\caption{数据分析：距离与辐射传感器电压的关系}
			\label{tab:tab5}
			\begin{tabular}{|c|c|}
				\hline
				距离/mm & 辐射传感器电压/mV \\
				\hline
				120 & 512.975 \\
				130 & 428.389 \\
				140 & 365.017 \\
				150 & 312.667 \\
				160 & 273.389 \\
				170 & 240.444 \\
				180 & 215.667 \\
				190 & 195.316 \\
				200 & 178.000 \\
				210 & 164.459 \\
				220 & 153.000 \\
				230 & 143.000 \\
				240 & 139.000 \\
				250 & 135.000 \\
				\hline
			\end{tabular}
		\end{table}
		
		
		\begin{figure}[htbp]
			\centering
			\includegraphics[width=0.6\textwidth]{Lab2_7Gra7.png}
			\caption{数据分析：不同辐射距离在$d$的辐射传感器输出$P$（球面波）}
			\label{fig:fig7}
		\end{figure}
		
		\item 从相关系数来看，拟合效果较好。接下来讨论分析这些拟合参数。因为辐射体并不是固定在标尺的零点位置，所以加入一个位移项偏置，功率项偏置用于消除环境的影响。但是，需要说明的是，这三个参数是拟合出来的，我在一定程度上怀疑它们是否是有效的（从程序运行结果上看，这样的拟合非常勉强；我尝试使用深度学习的办法对此做了一些探究，但并没有得到什么有用的结果，这里就不赘述了）。
		
		\item 现在有考虑讲义中提供的模型，考虑拟合函数$P=k\sin^2(\arctan\frac{R}{d-d_{bias}})+P_{bias}$拟合，其中$d_{bias}$可以考虑前一次拟合得到的结果，得到的拟合结果如\cref{fig:fig8}所示。
		
		\begin{figure}[htbp]
			\centering
			\includegraphics[width=0.6\textwidth]{Lab2_7Gra8.png}
			\caption{数据分析：不同辐射距离在$d$的辐射传感器输出$P$（平面辐射）}
			\label{fig:fig8}
		\end{figure}
		
		同样从相关系数看，拟合效果不错，但是从实际上去考虑拟合参数，效果不是很好。
		
		\item 对比这两个拟合结果，发现无论是参数k，还是两项偏置，差异都比较大，这说明实验数据与两个模型的符合情况都不太好。理论上讲，如果是从很近的距离开始测量，应该会发现使用平面辐射源的模型拟合结果更精确。
		
		\item 以下距离的定义哪个正确：（1）从辐射表面到传感器表面（红外透镜）的距离；（2）从辐射表面到传感器内部热电堆表面的距离。两者对测量误差的影响有多大？
		
		根据拟合结果，第二个距离的定义更准确，因为经过透镜只是用作滤去可见光的影响，而不是真正探测信号的元件；且辐射只有到达热电堆表面后，才会产生热量，才会有电压差输出。
		
		事实上，由于透镜的成像效果也许会改变等效的距离，因此在一定程度上上述结果可能都是错的。
	\end{enumerate}
	
	%
	\subsubsection{探讨测量不同物体表面的发射系数（Discussion）}
	根据测定的实验数据（\cref{tab:tab3})，我么可以做如下分析：
	
	发射系数是一个材料特性，指的是一个物体表面能够发射的辐射能与一个理想黑体在相同温度下发射的辐射能的比率。发射系数的数值范围从0到1，其中1代表完美的黑体，它可以吸收所有到达其表面的辐射能量，并且能以最大效率重新发射这些能量。
	
	\textbf{具体实验数据分析}
	
	\begin{enumerate}
		\item \textbf{光面}（传感器输出电压：285微伏）：
		\begin{itemize}
			\item 光面往往反射大部分入射辐射，吸收较少，因此其发射系数较低。这种表面在工程应用中常用于减少能量损失，例如在太阳能热水器的反射镜中。
		\end{itemize}
		
		\item \textbf{粗糙面}（传感器输出电压：313微伏）：
		\begin{itemize}
			\item 粗糙表面通过其不平整的表面特征增加了与辐射的相互作用，从而提高了吸收率。这种表面结构使得更多的能量被转换为热，然后以辐射的形式重新发射，导致其发射系数略高于光面。
		\end{itemize}
		
		\item \textbf{黑面}（传感器输出电压：1030微伏）：
		\begin{itemize}
			\item 黑色表面由于其特殊的物理和化学性质，能够吸收几乎所有入射的辐射能量，因此其发射系数接近1。这使得黑面具有非常高的能量转换效率，广泛应用于需要高吸收率的领域，如太阳能集热板和红外辐射测量。
		\end{itemize}
	\end{enumerate}
	
	
	% ---
	
	% 实验后思考题
	\clearpage
	\subsection{实验后思考题（Questions）}
	
	%思考题1
	\begin{question}
		为何当热辐射传感器太靠近辐射表面时，所测出的辐照度偏离距离平方反比规律？
	\end{question}
	具体原因已经在讲义中解释得非常清楚了，如\cref{fig:fig9}所示。
	
	\begin{figure}[htbp]
		\centering
		\includegraphics[width=0.7\textwidth]{Lab2_7Gra9.png}
		\caption{平面辐射源讲解}
		\label{fig:fig9}
	\end{figure}
	
	简言之，当距离靠的过近时，传感器能接收到的辐射体面积显著减小，随着距离的增大，能接收到的辐射体面积也增大。除此之外，我们还可以进一步考虑热辐射传感器太靠近辐射表面时，根据麦克斯韦方程组，电磁波在靠近辐射源的近场区域内，其电场和磁场并不遵循简单的距离平方反比规律，而是会呈现出更复杂的衰减特性。这是因为在近场区域内，电荷的加速变化所产生的电磁波不仅包含远场辐射，还包含近场感应和近场辐射。近场感应和近场辐射的衰减规律与距离平方反比规律不同，因此会导致传感器测得的辐照度偏离距离平方反比规律。	热辐射传感器本身也是一个辐射源，会向周围环境辐射电磁波。当传感器太靠近辐射表面时，传感器自身辐射的电磁波会与辐射表面的辐射电磁波混叠在一起，导致传感器测得的辐照度偏高。此外，传感器本身的结构和材料也会对测得的辐照度产生影响。
	
	% 思考题2
	\begin{question}
		SMTIR9902与SMTIR9902SIL各有什么特点？各适合于什么场景的非接触温度测量？为什么？
	\end{question}
	\begin{itemize}
		\item 根据查找到的资料，SMTIR9902 和 SMTIR9902SIL 是 Smartec 公司生产的红外传感器，专门用于非接触式温度测量。这两款传感器都具有高精度、高灵敏度（110 V/W）、优异的信噪比（约50KΩ）、快速的响应时间（40毫秒）和低成本的薄膜技术。它们适用于各种需要远距离、连续、非接触式的温度测量场合。
		
		SMTIR9902 包含一个内置的温度传感器，用于测量传感器自身的温度，非常适合于需要传感器自温补偿的应用场景，如环境温度控制和医疗器械。
		
		SMTIR9902SIL 在 SMTIR9902 的基础上加装了硅透镜，这使得它可以在更小的视角（7.27\degree）内测量辐射温度，因此非常适合于需要精确测量小区域或远距离对象的场景，例如在特定的工业应用中监测设备的热点或在医疗领域中精确测量人体的特定部位的温度。这种小视角和高精度的组合使得SMTIR9902SIL特别适用于精细和精密的温度测量任务。
		
		\item SMTIR9902：更适用于广泛的工业和商业应用，如HVAC系统的温度监控、电子设备的热管理、和食品加工的温度控制。该型号足以应对大部分标准的热测量需求，特别是在对象较大或温度分布较均匀的情况下。
		
		\item SMTIR9902SIL：适合于精细的科学研究，需要精确测量小区域或细节的场合，如材料科学中的热性质研究，或者在半导体制造过程中对芯片的热分析。 硅透镜提供的改进光学特性使得这款传感器更适用于对光学性能有较高要求的应用，例如在医疗成像和高精度工业过程控制中。
	\end{itemize}
	
	% 思考题3
	\begin{question}
		辐射传感器输出信号与辐射表面绝对温度4次方成线性关系，为何直线不过原点？Stefan-Boltzmann定律是否成立？为什么？
	\end{question}
	\begin{itemize}
		\item 事实上，我们已经结合具体实验数据和拟合结果作了分析，详见\textbf{数据分析}部分。
		下面进行详细的总结。
		
		\item  为什么直线不过原点？
		在实际测量中，辐射传感器输出信号与辐射体表面绝对温度的四次方的关系可能不完美地符合Stefan-Boltzmann定律。这可能有几个原因：传感器可能在测量时同时捕获来自周围环境的辐射（如室温辐射、仪器自身发出的热辐射等），这些背景辐射即使在目标辐射体温度为绝对零度时也可能存在，导致输出信号不为零。传感器和电子设备的内部噪声、偏置电流或电压也可能导致测量值在辐射体温度非常低时仍不为零。传感器的校准过程中可能包括了固有偏差的修正，使得在辐射体温度接近零的条件下仍有基线输出。 通常会对传感器进行校准以抵消传感器和电路的非理想特性，如零点偏移和非线性响应。 测量过程中的实验误差，如温度读数的不准确、传感器放置的位置、角度差异等，都可能影响最终的输出信号。物理限制如传感器的感光元件和材料特性也可能使得理论上的 \( T^4 \) 关系在低温区域有所偏离。
		
		\item Stefan-Boltzmann定律是否成立？
		Stefan-Boltzmann定律在理论上是成立的，并在实验中被广泛验证。如果辐射传感器输出信号与 \( T^4 \) 关系的直线不过原点，这通常不表示定律不成立，而是反映了实际应用中的各种复杂因素的影响。这些因素包括仪器的非理想行为、环境干扰、测量技术的限制等。
	\end{itemize}
	
	
	% 思考题4
	\begin{question}
		如何评价材料（如透光隔热膜）的防热辐射能力及其在窗户应用中的节能效果？
	\end{question}
	\begin{itemize}
		\item 材料的热性能评估
		\begin{itemize}
			\item \textbf{太阳能透过率}（Solar Transmittance）：这是指材料允许多少太阳能通过的比例，通常通过测定在特定波长范围内（如300nm至2500nm）的透光率来评估。
			\item \textbf{太阳能反射率}（Solar Reflectance）：这是指材料反射太阳能的能力，同样需要在特定的波长范围内进行测量。
			\item \textbf{太阳能吸收率}（Solar Absorptance）：这是指材料吸收太阳能的比例。理想的透光隔热膜应该有低的吸收率，以减少内部温度升高。
			\item \textbf{总太阳能透过系数}（Total Solar Energy Transmitted，g-value）：这是一个综合考量透过率、反射率和吸收率的指标，用于评估材料的总体热性能。
		\end{itemize}
		
		\item 实验方案设计
		
		为了准确评估透光隔热膜的防热辐射能力及其节能效果，可以设计以下实验：
		
		\begin{enumerate}
			\item 实验设备和材料
			\begin{itemize}
				\item \textbf{透光隔热膜样品}
				\item \textbf{光源设备}：模拟太阳光的光源
				\item \textbf{温度传感器}：用于测量膜两侧的温度变化
				\item \textbf{紫外线-可见光-近红外光谱仪}：用于测定透过率、反射率和吸收率
				\item \textbf{热像仪}：用于监测热辐射的分布和变化
			\end{itemize}
			
			\item 实验步骤
			\begin{enumerate}
				\item \textbf{样品准备}：将透光隔热膜安装在一个标准窗户模型上。
				\item \textbf{环境设置}：在模拟环境中，使用光源设备模拟不同强度和角度的太阳光。
				\item \textbf{数据采集}：使用温度传感器监测膜的两侧以及室内外的温度差异，使用光谱仪测量太阳能透过率、反射率和吸收率。
				\item \textbf{数据分析}：分析数据以评估透光隔热膜的热性能和节能效果。关注在不同条件下（如不同季节、不同时间段）的节能效果。
				\item \textbf{优化与调整}：根据初步实验结果，对材料或实验设置进行优化，如调整膜的厚度或改变安装角度等，以获得更好的节能效果。
			\end{enumerate} 
		\end{enumerate}
		
		\item 结果评估和报告
		\begin{itemize}
			\item 准备详细的实验报告，包括实验条件、过程、数据分析和结论。
			\item 对比不同材料和不同安装方式的节能效果，为实际应用提供科学依据。
		\end{itemize}
		
	\end{itemize}
	
	
	% 思考题5
	\begin{question}
		如何在（变化的）大气环境温度下测量各热辐射源的温度？
	\end{question}
	实验方案：红外热像仪在变化环境下的温度测量
	\begin{itemize}
		\item 实验目的
		
		验证红外热像仪在不同环境温度下测量不同热辐射源温度的准确性和可靠性。
		
		\item 实验材料
		\begin{enumerate}
			\item 红外热像仪
			\item 多个不同材料的热辐射源（如金属块、塑料块）
			\item 可调节的环境温度设备（如恒温箱）
			\item 标准温度计或热电偶（用于校准和比较）
		\end{enumerate}
		
		\item 实验步骤
		\begin{enumerate}
			\item \textbf{设备校准}：确保红外热像仪已经根据制造商的指南进行了校准。
			\item \textbf{设置环境}：将恒温箱设置在一个初始温度（例如20°C）。允许设备在此环境温度下稳定运行至少30分钟。
			\item \textbf{温度设定}：将不同材料的热辐射源预热到不同的目标温度（例如50°C、70°C和90°C）。
			\item \textbf{测量记录}：
			\begin{itemize}
				\item 使用红外热像仪从同一距离测量每个热源的表面温度。
				\item 记录热像仪显示的温度值。
				\item 同时使用标准温度计或热电偶测量并记录实际温度。
			\end{itemize}
			\item \textbf{环境温度变化}：调整恒温箱的温度至不同水平（例如10°C、30°C），重复步骤3和4。
			\item \textbf{数据分析}：分析红外热像仪和标准温度计的温度读数差异，考察不同环境温度对测量结果的影响。
		\end{enumerate}
		
		\item 数据处理
		\begin{itemize}
			\item 对比红外热像仪和标准温度计的读数，计算误差。
			\item 分析误差随环境温度变化的趋势。
		\end{itemize}
		
		\item 预期结果
		\begin{itemize}
			\item 验证红外热像仪在不同环境温度下的测量精确度。
			\item 确定最适合的环境温度范围，以便在实际应用中使用红外热像仪获取最准确的结果。
		\end{itemize}
	\end{itemize}
	
	% ---
	
	
	% 结语部分
	\clearpage
	
	% 小标题
	\section{Lab2-7 热辐射实验（The Experiments on Thermal Radiation）\\ \quad\heiti 结语（The End）}
	% ---
	
	% 总结、杂谈与致谢
	\subsection{总结、杂谈与致谢（Summary, Thoughts \& Thanks）}
	\begin{enumerate}
		\item \textbf{数据处理代码参考了戴鹏辉22344016同学编写的python脚本并作了适当的改进与调整}，在此特别表示感谢。同时，也感谢老师在实验过程中的帮助。
		\item 感谢老师能阅读这篇还有很多不足的实验报告，希望老师能指出做得不好的地方；祝老师身体健康、生活幸福、工作顺利！
	\end{enumerate}
	% ---
	
	% 参考文献
	\subsection{参考文献（References）}
	[1] 维基百科 https://zh.wikipedia.org
	
	[2] 沈韩.基础物理实验.——北京：科学出版社，2015.2 ISBN：978-7-03-043311-4
	
	% ---
	
	% 附件
	\subsection{附件（Attachment）}
	%实验台桌面整理如%\cref{}所示。
	
	实验报告个人签名如\cref{fig:name}。
	
	\begin{figure}[htbp]
		\centering
		\includegraphics[width=0.7\textwidth]{name.png}
		\caption{个人签名}
		\label{fig:name}
	\end{figure}
	
	% ---
	
	相关代码已上传至Github。
	
	
	
\end{document}